\chapter{pandas による表形式データ解析}
\label{chap:pandas}

pandas\index{pandas@pandas|textbf} は、行と列からなる表形式データを扱うライブラリである。観測ログ、イベント一覧、較正表、集計結果のように、各行が一つの記録、各列が一つの変数に対応するデータを \code{DataFrame} で表現できる。本章では、列名と型の確認、条件抽出、欠損値処理、集計、表変形、入出力を扱う。

\section{DataFrame の構造と列指向処理}

\subsection{Series と DataFrame}

pandas の基本構造は \code{Series}\index{series@Series|textbf} と \code{DataFrame}\index{dataframe@DataFrame|textbf} である。\code{Series} はラベル付き 1 次元配列、\code{DataFrame} は列名付き 2 次元表である。NumPy 配列が配列全体で一つの \code{dtype} を持つのに対し、DataFrame は列ごとに異なる型と列名を保持できる。

\begin{lstlisting}[language=Python]
import pandas as pd

events = pd.DataFrame(
    {
        "detector": ["A", "A", "B", "B"],
        "time_ns": [1032, 1188, 1055, 1244],
        "signal": [12.4, 10.1, 14.8, 9.7],
        "valid": [True, True, False, True],
    }
)

print(type(events).__name__)
print(events)
\end{lstlisting}

このコードを実行すると、次のような DataFrame が表示される。

\begin{lstlisting}[numbers=none, xleftmargin=2\zw]
DataFrame
  detector  time_ns  signal  valid
0        A     1032    12.4   True
1        A     1188    10.1   True
2        B     1055    14.8  False
3        B     1244     9.7   True
\end{lstlisting}

この表では、行番号 0, 1, 2, 3 は index、\code{"detector"} や \code{"signal"} は列名である。index は行を識別するラベルであり、列名は各列が表す量を示す。値だけでなく、列名、型、単位、欠損値の有無を合わせて確認する。

DataFrame の利点は、列名によって意味を示しながら、列単位の処理を記述できることにある。列の位置を整数で保持する必要がなくなり、ファイル読み込み、列変換、フィルタ、集計、書き出しまでを一つの表に対する処理として扱える。

DataFrame の列選択は辞書のキーによる値の参照、行の条件抽出はブール配列による NumPy 配列の抽出に対応する。ただし、DataFrame では index と列名のラベルも演算時の対応付けに使われる。

\subsection{行例・形状・列型・欠損値の確認}
\index{pandas!describe@\code{describe()}}

DataFrame の作成または読み込みの直後には、\code{head()}、\code{shape}、\code{dtypes}、\code{info()} で行の例、形状、列型、欠損値数を確認する。

\begin{lstlisting}[language=Python]
print(events.head())
print(events.shape)
print(events.dtypes)
\end{lstlisting}

出力は次の通りである。

\begin{lstlisting}[numbers=none, xleftmargin=2\zw]
  detector  time_ns  signal  valid
0        A     1032    12.4   True
1        A     1188    10.1   True
2        B     1055    14.8  False
3        B     1244     9.7   True
(4, 4)
detector        str
time_ns       int64
signal      float64
valid          bool
dtype: object
\end{lstlisting}

\code{shape} は行数と列数、\code{dtypes} は各列の型である。\code{object} 型は任意の Python オブジェクトを格納できるため、その表示だけでは列が文字列だけで構成されるか、複数の型が混在するかを判定できない。数値として演算する列が \code{object} 型なら、実際の値と変換条件を確認する。

上の型表示は pandas 3 系の例である。文字列だけの列は既定で \code{str} と推定され、pandas 2 系の既定では \code{object} となる~\cite{pandas_migration3}。表示の違いと、数値型への変換の要否は分けて確認する。

\subsection{index と列名の役割}
\index{pandas!index@index}

pandas では、整数位置とラベルを区別して行と列を選択する。

\begin{lstlisting}[language=Python]
print(events["signal"])
print(events.loc[1:2, ["detector", "signal"]])
print(events.iloc[1:3, 0:2])
\end{lstlisting}

\code{events["signal"]} は列名による選択、\code{loc}\index{loc@\code{loc}|textbf} はラベルによる選択、\code{iloc}\index{iloc@\code{iloc}|textbf} は整数位置による選択である。pandas の indexing ガイド~\cite{pandas_indexing} でも、選択対象に応じて \code{[]}、\code{loc}、\code{iloc}、\code{at}、\code{iat} を使い分けている。解析コードでは、列名・ラベル・整数位置のどれを指定しているかを明確にする。

\subsection{Python の行反復と列演算の実装差}

pandas の列演算は、要素ごとに Python 関数を呼び出すループを避け、配列単位の実装へ処理を渡せる。実行時間の差はデータ型と処理内容に依存するため、対象データを用いて列演算と行反復を測定する。

ただし、すべての処理が pandas の列演算に対応するわけではない。行間の近傍探索や状態遷移のように、前後関係または専用のデータ構造を必要とする処理では、NumPy、探索アルゴリズム、または目的に合わせたデータ構造を検討する。処理を独立した列の変換・抽出・集計として表せるかどうかが、pandas を選択する基準となる。

\section{抽出・変換・集計・入出力}

\subsection{列選択と行抽出}

列選択と条件による行抽出を組み合わせる例を示す。

\begin{lstlisting}[language=Python]
selected = events[["detector", "signal"]]
valid_events = events.loc[events["valid"], :]
strong = events.loc[events["signal"] > 11.0, ["detector", "signal"]]

print(selected)
print(valid_events)
print(strong)
\end{lstlisting}

各変数に残る行と列は次のとおりである。

\begin{lstlisting}[numbers=none, xleftmargin=2\zw]
  detector  signal
0        A    12.4
1        A    10.1
2        B    14.8
3        B     9.7
  detector  time_ns  signal  valid
0        A     1032    12.4   True
1        A     1188    10.1   True
3        B     1244     9.7   True
  detector  signal
0        A    12.4
2        B    14.8
\end{lstlisting}

ブール配列で行を絞る操作は NumPy と似ているが、pandas では列名を介して意味を保ったまま書ける。\code{events["signal"] > 11.0} の結果は行ごとの真偽値であり、それを \code{loc} に渡すと条件に合う行だけが残る。

\code{loc} は行ラベルや列名で選択し、ラベルのスライスでは終点も含む。\code{iloc} は 0 から始まる位置で選択し、Python の通常のスライスと同じく終点を含まない。また、Series を列へ代入すると、通常は位置ではなく index のラベルで対応付けられる。

\subsection{列の追加}

元の列から比、差、判定結果を求める場合は、列全体を入力とする式を記述する。

\begin{lstlisting}[language=Python]
events["signal_ratio"] = events["signal"] / events["signal"].mean()
events["is_late"] = events["time_ns"] > 1100

print(events)
\end{lstlisting}

列を追加すると、元の表へ新しい列が右側へ増える。出力は次のようになる。

\begin{lstlisting}[numbers=none, xleftmargin=2\zw]
  detector  time_ns  signal  valid  signal_ratio  is_late
0        A     1032    12.4   True      1.055319    False
1        A     1188    10.1   True      0.859574     True
2        B     1055    14.8  False      1.259574    False
3        B     1244     9.7   True      0.825532     True
\end{lstlisting}

\code{assign()} を使うと、途中の DataFrame を明示しながら複数列をまとめて追加できる。

\begin{lstlisting}[language=Python]
derived = (
    events
    .assign(
        dt_ns=lambda df: df["time_ns"] - df["time_ns"].min(),
        quality=lambda df: pd.Series(["high", "mid", "low", "mid"])
    )
)
\end{lstlisting}

\code{assign()} は新しい DataFrame を返すため、メソッドチェーンの途中で列を追加できる。一方、\code{events["new\_col"] = ...} は既存の DataFrame に列を直接追加する。元の表を変更するか、新しい表を返すかに応じて選択する。

\subsection{欠損値と型}
\index{pandas!isna@\code{isna()}}
\index{けっそんち@欠損値|textbf}

現実のデータでは、空欄や壊れた値が入り込む。pandas ではこれを \code{NaN} や \code{NA} として扱うことが多い。

\begin{lstlisting}[language=Python]
df = pd.DataFrame(
    {
        "id": [1, 2, 3, 4],
        "signal": [12.4, None, 14.8, 9.7],
        "comment": ["ok", "bad", None, "ok"],
    }
)

print(df.isna())
print(df.dropna(subset=["signal"]))
print(df["signal"].fillna(df["signal"].mean()))
\end{lstlisting}

\code{isna()} の真偽表、\code{dropna()}\index{pandas!dropna@\code{dropna()}} の結果、\code{fillna()}\index{pandas!fillna@\code{fillna()}} 後の列を順に表示すると、欠損処理の違いが見えやすい。

\begin{lstlisting}[numbers=none, xleftmargin=2\zw]
      id  signal  comment
0  False   False    False
1  False    True    False
2  False   False     True
3  False   False    False
   id  signal comment
0   1    12.4      ok
2   3    14.8     NaN
3   4     9.7      ok
0    12.400000
1    12.300000
2    14.800000
3     9.700000
Name: signal, dtype: float64
\end{lstlisting}

欠損値は \code{==} で判定しない。\code{isna()}、\code{notna()}、\code{dropna()}、\code{fillna()} を使う。また、整数列に欠損が入ると \code{float64} に変わることがあるので、必要なら pandas の nullable integer 型である \code{Int64} を使う。

\begin{lstlisting}[language=Python]
df["id"] = df["id"].astype("Int64")
\end{lstlisting}

\subsection{要約統計と \code{groupby}}

\code{describe()} は数値列の件数、平均、標準偏差、分位点、最小値、最大値を返す。\code{groupby()}\index{groupby@\code{groupby}|textbf} は、カテゴリごとに分割した行へ同じ集計を適用する。

\begin{lstlisting}[language=Python]
print(events.describe())

summary = (
    events
    .groupby("detector")
    .agg(
        count=("signal", "size"),
        mean_signal=("signal", "mean"),
        std_signal=("signal", "std"),
        late_fraction=("is_late", "mean"),
    )
)

print(summary)
\end{lstlisting}

\code{describe()} は全体の要約、\code{groupby()} の結果は検出器ごとの集計である。出力は次の通りである。

\begin{lstlisting}[numbers=none, xleftmargin=2\zw]
           time_ns     signal  signal_ratio
count     4.000000   4.000000      4.000000
mean   1129.750000  11.750000      1.000000
std     102.613758   2.355844      0.200497
min    1032.000000   9.700000      0.825532
25%    1049.250000  10.000000      0.851064
50%    1121.500000  11.250000      0.957447
75%    1202.000000  13.000000      1.106383
max    1244.000000  14.800000      1.259574
          count  mean_signal  std_signal  late_fraction
detector
A             2        11.25    1.626346            0.5
B             2        12.25    3.606245            0.5
\end{lstlisting}

\code{groupby()} は、検出器、日、run などのカテゴリごとに同じ集計を適用する。\code{agg()} の名前付き集計では、出力列の名前、入力列、集計関数を対応付けて記述できる。上の数値は、この節の \code{events} を pandas で計算した出力例である。pandas のバージョンによって空白や型の表示形式は変わる場合があるが、数値は同じ入力と計算定義に対応していなければならない。

欠損値を含む列の平均や標準偏差は、既定では欠損値を除いて計算する。\code{groupby().size()} は行数を、列に対する \code{count()} は非欠損値の個数を数える。\code{std()} の既定は \code{ddof=1} であり、有効値が 1 個だけの群では標準偏差は \code{NaN} になる。未知の文字列は必ずしも欠損値として認識されない。\code{pd.to\_numeric(..., errors="coerce")} で変換できなかった値を欠損にする場合は、その件数と元の値を記録する。平均値による補完は分散などを変えるため、解析上の根拠なしに適用しない。

\subsection{表形式の変換}
\index{pandas!merge@\code{merge()}}
\index{pandas!concat@\code{concat()}}
\index{pandas!pivot_table@\code{pivot\_table()}}

Getting started tutorials~\cite{pandas_tutorials} でも、表を縦長から横長へ変える操作や、複数表を結合する操作が扱われている。解析では、観測値と較正値を結び付けたり、日別表をまとめたりする場面で用いる。

\begin{lstlisting}[language=Python]
gain = pd.DataFrame(
    {
        "detector": ["A", "B"],
        "gain": [1.02, 0.97],
    }
)

merged = events.merge(gain, on="detector", how="left", validate="many_to_one")
assert len(merged) == len(events)
stacked = pd.concat([events.iloc[:2], events.iloc[2:]], ignore_index=True)

long_df = pd.DataFrame(
    {
        "detector": ["A", "A", "B", "B"],
        "kind": ["raw", "cal", "raw", "cal"],
        "value": [12.4, 12.7, 14.8, 14.4],
    }
)
wide = long_df.pivot(index="detector", columns="kind", values="value")

print(merged[["detector", "signal", "gain"]])
print(wide)
\end{lstlisting}

次の出力は、\code{merge()} による較正係数列の追加と、\code{pivot()} による種別 \code{kind} の値から列への変換を示す。

\begin{lstlisting}[numbers=none, xleftmargin=2\zw]
  detector  signal  gain
0        A    12.4  1.02
1        A    10.1  1.02
2        B    14.8  0.97
3        B     9.7  0.97
kind       cal   raw
detector
A         12.7  12.4
B         14.4  14.8
\end{lstlisting}

\code{merge()} はキー列に基づいて表を結合し、\code{concat()} は表を行方向または列方向に連結する。\code{pivot()} と \code{pivot\_table()} は縦長データを横長形式へ変換し、\code{melt()} は横長データを縦長形式へ変換する。

結合先のキーが重複していると、一つの観測行に複数の行が対応し、結合後の行数が増える。上の \code{validate="many\_to\_one"} は較正表側のキーの一意性を検査する。結合できなかったキーでは較正値が欠損するため、行数に加えて欠損件数も確認する。\code{pivot()} は同じ行キーと列キーの組合せが重複するとエラーになる。重複を集計する場合は \code{pivot\_table(..., aggfunc="mean")} のように集計方法を明示する。

\subsection{時刻、文字列、カテゴリ}
\index{pandas!datetime@datetime}
\index{pandas!category@category}

時刻や文字列は、そのままでは解析しにくい。pandas には時刻用の \code{dt} アクセサ、文字列用の \code{str} アクセサ、カテゴリ型がある。

\begin{lstlisting}[language=Python]
log = pd.DataFrame(
    {
        "timestamp": ["2026-04-01 12:00:00", "2026-04-01 12:00:05", "2026-04-02 09:30:00"],
        "run_name": ["run_A_001", "run_A_002", "test_B_001"],
        "detector": ["A", "A", "B"],
    }
)

# 元のタイムゾーンなし文字列が UTC を表すと確認した場合
log["timestamp"] = pd.to_datetime(log["timestamp"], utc=True)
log["date"] = log["timestamp"].dt.date
log["prefix"] = log["run_name"].str.extract(r"^(run_[A-Z])")
log["detector"] = log["detector"].astype("category")

print(log[["timestamp", "date", "prefix", "detector"]])
\end{lstlisting}

時刻列と文字列列を分解すると、日付や接頭辞を独立した列として扱える。出力は次の通りである。

\begin{lstlisting}[numbers=none, xleftmargin=2\zw]
                  timestamp        date prefix detector
0 2026-04-01 12:00:00+00:00  2026-04-01  run_A        A
1 2026-04-01 12:00:05+00:00  2026-04-01  run_A        A
2 2026-04-02 09:30:00+00:00  2026-04-02    NaN        B
\end{lstlisting}

\code{to\_datetime()} で日時型へ変換すると、時刻差、日付成分、時刻順を日時の演算として扱える。文字列列には \code{str.contains()}、\code{str.extract()}、\code{str.strip()} などのベクトル化メソッドがある。カテゴリ型は、有限個のカテゴリ値と整数コードを分けて保持し、カテゴリの順序も定義できる。

\subsection{入出力}
\index{parquet@Parquet|textbf}
\index{pandas!parquet@Parquet}

pandas は、CSV、Parquet、圧縮テキストなどを DataFrame として読み書きできる。本節では、ファイル形式ごとの型と圧縮方法の指定を扱う。

\index{pandas!read\_csv@\code{read\_csv()}|textbf}
\begin{lstlisting}[language=Python]
df = pd.read_csv(
    "events.csv.gz",
    compression="gzip",
    usecols=["detector", "time_ns", "signal"],
    dtype={"detector": "category", "time_ns": "int64", "signal": "float64"},
)

df.to_parquet("events.parquet")
small = pd.read_parquet("events.parquet", columns=["detector", "signal"])
\end{lstlisting}

この例は、ヘッダ付きのカンマ区切り CSV を gzip 圧縮した \code{events.csv.gz} を想定している。第 1 引数は入力パス、\code{compression="gzip"} は圧縮方式、\code{usecols} は読み込む列、\code{dtype} は列ごとの型である。\code{to\_parquet()} は表全体を Parquet へ保存し、\code{read\_parquet(..., columns=...)} は指定した 2 列だけを読み込む。

\code{read\_csv()} では、必要な列、列の型、コメント行の扱いを入力時に指定できる。次の例は、gzip 圧縮された空白区切りファイルから、必要な列だけを指定した型で読み込む。

\begin{lstlisting}[language=Python]
df = pd.read_csv(
    path,
    compression="gzip",
    sep=r"\s+",
    comment="#",
    header=None,
    usecols=[0, 1, 2, 3],
    names=["id", "date", "time", "value"],
    dtype={"id": "int64", "value": "float64"},
)
\end{lstlisting}

必要な列だけを読み込むと、読み込み時間とメモリ消費を減らせる場合がある。さらに \code{dtype=} を指定すれば、読み込み後の型変換も減らせる。列数や行数が大きいデータでは、必要な列を入力段階で選ぶことが後続処理の負荷軽減につながる。

2 つ目の例は、ヘッダのない空白区切りファイルを想定している。\verb|sep=r"\s+"| は 1 個以上の空白を区切りとし、\verb|comment="#"| は \code{\#} から始まる行の残りをコメントとして扱う。\code{header=None} は列名行がないことを示し、\code{names} が読み込んだ 4 列へ名前を付ける。\code{usecols=[0, 1, 2, 3]} の番号は入力ファイル中の列位置であり、0 から数える。

大規模なテキストファイルの反復読み込みが処理時間を支配する場合は、一度整形した表を列型と変換条件とともに Parquet へ保存する。\code{read\_parquet()} では必要な列だけを選択できるため、反復解析の読み込み時間とメモリ使用量を削減できる場合がある。効果は入力容量、列数、ストレージによって変わるため、変換前後を測定する。

Parquet はファイル形式の名前であり、\code{pyarrow} はその読み書きを担うライブラリである。\code{to\_parquet()} や \code{read\_parquet()} がライブラリ不足を示すエラーで失敗した場合は、Parquet ファイルの破損ではなく、\code{pyarrow} などの対応ライブラリが環境に入っていない可能性がある。

列同士の演算をそのまま式として書ける点も pandas の強みである。

\begin{lstlisting}[language=Python]
df["ratio"] = df["value"] / df["value"].mean()
selected = df[df["ratio"] > 1.5]
\end{lstlisting}

行ごとの \code{for} ループを導入する前に、同じ処理を列演算、集計、結合として表せるかを検討する。

Parquet の例には、たとえば \code{python -m pip install pyarrow} で導入できるエンジンが必要である。後の \code{to\_latex()} には Jinja2 が必要なので、\code{python -m pip install jinja2} で導入する。uv のプロジェクトではそれぞれ \code{uv add} を用いる。

\subsection{大規模 CSV の分割集計}
\index{pandas!chunksize@\code{chunksize}}
\index{ぶんかつよみこみ@分割読み込み|textbf}

pandas の DataFrame は主としてメモリ上で処理されるが、\code{read\_csv()} の \code{chunksize} を使えば、CSV を一定行数ずつ読みながら集計できる。例として、次の \code{events.csv} を考える。

\begin{lstlisting}[numbers=none, xleftmargin=2\zw, title={events.csv}]
detector,signal
A,12.4
B,14.8
A,10.1
B,9.7
A,13.0
\end{lstlisting}

検出器ごとの件数と信号値の合計は、次のように 2 行ずつ読んでも計算できる。

\begin{lstlisting}[language=Python]
counts = pd.Series(dtype="int64")
sums = pd.Series(dtype="float64")

for chunk in pd.read_csv(
    "events.csv",
    usecols=["detector", "signal"],
    dtype={"detector": "string", "signal": "float64"},
    chunksize=2,
):
    grouped = chunk.groupby("detector")["signal"]
    counts = counts.add(grouped.count(), fill_value=0)
    sums = sums.add(grouped.sum(), fill_value=0)

summary = pd.DataFrame({"count": counts, "mean": sums / counts})
print(summary)
\end{lstlisting}

実行結果は次の通りである。

\begin{lstlisting}[numbers=none, xleftmargin=2\zw]
          count       mean
detector
A           3.0  11.833333
B           2.0  12.250000
\end{lstlisting}

\code{chunksize=2} は、反復ごとに最大 2 行の DataFrame を返す指定である。実データでは 10 万行など、利用可能なメモリと処理時間を見ながら値を決める。\code{Series.add(..., fill\_value=0)} は、あるチャンクに片方の検出器しか含まれない場合にも、存在しない集計値を 0 として足す。平均値をチャンク平均の単純平均として求めると、チャンクごとの件数が違う場合に誤るため、全体の合計を全体の件数で割っている。

分割処理が適するのは、総和、件数、最小値、最大値のようにチャンクごとの結果を統合できる計算である。全行を同時に並べ替える処理、全行どうしを結合する処理、全データに依存する順位計算では、単純な分割だけでは同じ結果にならない。その場合は、必要列を減らす、Parquet を使う、次章で説明する Polars の遅延実行やデータベースを使うといった別の方法も検討する。

\subsection{表の LaTeX 出力}
\index{pandas!to_latex@\code{to\_latex()}}

\code{to\_latex()} で集計結果から表を生成すれば、DataFrame の値と LaTeX 表の値を同じ処理から出力できる。

\begin{lstlisting}[language=Python]
from pathlib import Path

Path("tables").mkdir(parents=True, exist_ok=True)
tex_table = summary.reset_index().to_latex(
    index=False,
    caption="検出器ごとの平均信号",
    escape=True,
    label="tab:detector-summary",
    float_format="%.3f",
)

with open("tables/detector_summary.tex", "w", encoding="utf-8") as f:
    f.write(tex_table)
\end{lstlisting}

ここで \code{index=False} は左端の連番インデックスを出力しない指定であり、\code{reset\_index()} は groupby 結果の index を通常の列へ戻す操作である。ここでは \code{escape=True} によってセル中の TeX の特殊文字をエスケープする。pandas 2.0 以降では既定値に頼るとエスケープされない場合がある~\cite{pandas_latex}。検証済みの TeX 記法を意図的に埋め込む場合に限り、入力内容を限定したうえで \code{escape=False} を指定する。

\code{groupby()}、\code{agg()}、\code{to\_latex()} を順に適用し、保存したファイルをレポートから \verb|\input| すれば、集計結果をレポート用の表へ直接反映できる。この構成では、集計条件を変更した場合も同じ処理から表を再生成でき、転記による値の不一致を避けられる。

\subsection{\code{plot()} による簡易確認}

pandas の \code{plot()} は、DataFrame または Series の内容を Matplotlib へ渡す簡易作図インターフェースである。

\begin{lstlisting}[language=Python]
summary.plot(kind="bar", y="mean")
\end{lstlisting}

\code{DataFrame.plot()} は表から概形を確認する場合に用いる。提出図や論文用図では、軸、凡例、余白、保存形式を明示するため、Matplotlib の \code{Figure} と \code{Axes} を直接扱う。本書では第~\ref{sec:matplotlib-object}~節以降の方法を基準とする。

\section{コピー・型・行処理に関する注意点}

\subsection{部分表の変更と \code{loc}}
\index{settingwithcopy@SettingWithCopy|textbf}

DataFrame の一部を切り出した後の代入では、元の DataFrame を更新するのか、独立した表を作るのかを明示する必要がある。

\begin{lstlisting}[language=Python]
subset = events[events["signal"] > 11.0]
subset["flag"] = True
\end{lstlisting}

pandas 2 系の従来の動作では、この書き方で \code{SettingWithCopyWarning} が出ることがある。独立した部分表が必要なら抽出後に \code{.copy()} を付け、意図を明示する。pandas 3 系は Copy-on-Write が標準であり、部分表の変更は元の表へ反映されない~\cite{pandas_cow}。元の表の条件に合う行へ代入する場合は、次のように行条件と列名を \code{loc} の一つの添字指定へまとめる。

\begin{lstlisting}[language=Python]
events.loc[events["signal"] > 11.0, "flag"] = True
\end{lstlisting}

\subsection{\code{object} 型のまま計算しない}
\index{object dtype@\code{object} 型|textbf}

数字からなる文字列を含む列でも、空白や数字以外の記号が混在すると \code{object} 型で読み込まれる場合がある。そのまま演算すると、文字列連結または型に関する例外が生じる。入力仕様が決まっている列は \code{read\_csv()} の \code{dtype=} で型を指定し、読み込み後の \code{dtypes} と変換できなかった値を確認する。

\subsection{行ループを適用する前の検討事項}

pandas の DataFrame では、同じ処理でも列演算と行単位の Python 関数で実行時間が大きく異なることがある。特に、\code{iterrows()}、\code{apply(axis=1)}、並列化の使い分けは、処理内容と測定結果に基づいて判断する必要がある。

\subsubsection{\code{iterrows()} による行ループの負担}
\index{pandas!iterrows@\code{iterrows()}}
\index{pandas!itertuples@\code{itertuples()}}

DataFrame の行を一つずつ取り出す処理は記述しやすいが、列単位の演算よりオーバーヘッドが大きくなる場合がある。

\begin{lstlisting}[language=Python]
df = pd.DataFrame({"value": [1.0, 2.0, 3.0]})

def heavy_computation(val):
    return val ** 2

results = []
for index, row in df.iterrows():
    results.append(heavy_computation(row["value"]))
df["value_sq"] = results
\end{lstlisting}

\code{iterrows()} は各行を \code{Series} として生成するため、元の列型が行単位の Series に保たれない場合があり、行ごとのオブジェクト生成も生じる。行反復が必要なら \code{itertuples()} との処理時間を測定する。列ごとに独立して計算できる処理には、列演算または既存のベクトル化メソッドを用いる。

\subsubsection{\code{apply(axis=1)} と列演算の比較}

行ごとに処理を適用したい場合、pandas が用意している \code{apply()} メソッドを使うと、コードを DataFrame の処理として記述できる。\code{axis=1} は「各行へ関数を適用する」という意味であり、\code{axis=0} なら各列に適用する。

\begin{lstlisting}[language=Python]
df["value_sq"] = df.apply(lambda row: heavy_computation(row["value"]), axis=1)
\end{lstlisting}

ただし、\code{axis=1} を付けた \code{apply()} も各行に対して Python 関数を呼ぶため、列演算とは処理方式が異なる。単に二乗する場合は、次の列演算が演算対象と式を直接表す。実行時間は対象データを用いて測定する。

\begin{lstlisting}[language=Python]
df["value_sq"] = df["value"] ** 2
\end{lstlisting}

\code{apply()} は、既存の列演算や NumPy 関数では表しにくい処理に Python 関数を適用する手段である。列演算、\code{apply()}、並列化の候補について、同じ入力と同じ出力条件で処理時間を測定して選ぶ。

以下の \code{pandarallel} は追加ライブラリであり、使用する環境に \code{python -m pip install pandarallel} などで導入する。二乗程度の軽い処理では並列化の起動・転送コストが支配的になり得るため、この例自体が高速化の推奨例というわけではない。

\subsubsection{\code{pandarallel} による行処理の並列化}

通常の \code{apply()} は、Python 関数を行ごとに順番に呼び出す。本当に計算が重く、列演算や既存の NumPy 関数へ置き換えられない場合に、マルチプロセスによる並列化が選択肢に入る。

\begin{lstlisting}[language=Python]
from pandarallel import pandarallel

pandarallel.initialize(progress_bar=True)
df["value_sq"] = df.parallel_apply(
    lambda row: heavy_computation(row["value"]),
    axis=1,
)
\end{lstlisting}

図~\ref{fig:pandas-benchmark} は、\code{iterrows()}、\code{apply()}、\code{pandarallel} の処理時間に生じ得る違いを模式的に示す。実測値ではなく、処理方式を比較するときに確認すべきオーバーヘッドの違いを表した図である。

\begin{figure}[htbp]
\centering
\includegraphics[width=0.7\linewidth]{pandas_benchmark.pdf}
\caption{同じ行単位処理に対する Python のループ、\code{apply()}、\code{pandarallel} の模式的な比較。棒の高さは実測値ではなく、入力サイズや 1 行当たりの計算量によって順位も変わり得る。}
\label{fig:pandas-benchmark}
\end{figure}

\code{pandarallel} は、DataFrame を複数プロセスへ分割し、\code{parallel\_apply()} に渡した関数を各部分へ適用する。プロセスの生成、データの直列化、プロセス間転送の時間も加わるため、対象処理について逐次実行との実行時間とメモリ使用量を測定する。

\code{pandarallel} は複数の Python プロセスを起動するため、OS によってプロセスの開始方法や必要な設定が異なる。Windows で本書の Linux 系コマンドと実行条件をそろえる場合は、第~\ref{chap:shell}~章で説明した WSL を利用できる。

\code{pandarallel} は DataFrame を複数プロセスへ分配するため、プロセスの起動、データの受け渡し、結果の結合にも時間がかかる。単純な加算のように 1 行当たりの処理が短い場合は、この付加時間が計算時間を上回り、逐次処理より遅くなることがある。並列化の採否は、同じ入力に対する逐次処理と並列処理を測定して判断する。

\subsection{欠損値と比較演算を混同しない}

\code{NaN == NaN} は \code{False} であるため、欠損判定には \code{==} ではなく \code{isna()} を用いる。\code{None}、\code{NaN}、\code{NA} の使用可否と表示は列型に依存するので、\code{dtype} と欠損値数を確認する。

\section{観測表の集計と可視化への受け渡し}

観測データの解析では、pandas はしばしば次の流れの中心になる。

\begin{enumerate}
\item \code{read\_csv()} によるテキスト・CSV の読み込み
\item \code{head()} と \code{dtypes} による表構造の確認
\item 必要な列の選択と、欠損値・異常値の整理
\item \code{groupby()} と \code{merge()} による集計・結合
\item Matplotlib への受け渡しと作図
\item \code{to\_latex()} による集計表の書き出し
\end{enumerate}

この流れでは、列名と列型を保ったまま抽出、集計、結合、保存を記述している。観測条件、検出器 ID、run 名、欠損フラグのように型の異なる列を一つの記録として扱う場合は、列ごとに型を持つ DataFrame を利用できる。

一方、CSV の読み込みや集計が実測した待ち時間の主要部分を占める場合は、次章で扱う Polars も候補となる。行数だけで切り替えを決めず、必要な機能、メモリ使用量、処理時間を同じ入力で比較する。

また、pandas の多くの操作は、DataFrame が 1 台の計算機のメモリに収まることを前提とする。ただし、\code{read\_csv(chunksize=...)} を使えば、総和や件数のように部分結果を統合できる処理は全件を同時に保持せずに実行できる。必要な表がメモリに収まらない場合は、列と行を早い段階で絞る、分割集計を使う、Parquet に変換する、次章の Polars やデータベースを使う、という順に処理の性質に合う方法を選ぶ。

基本操作の公式資料として Getting started tutorials~\cite{pandas_tutorials} と 10 minutes to pandas~\cite{pandas_10min} がある。列選択と代入は Indexing~\cite{pandas_indexing}、欠損値処理は Missing data~\cite{pandas_missing}、表の結合は Merge, join, concatenate and compare~\cite{pandas_merging}、時刻列は Time series / date functionality~\cite{pandas_timeseries} に詳しい。本章は、これらの機能を解析手順に沿って再構成した。

\section{章末課題：検出器別の表集計}

次の CSV を \code{events.csv} として保存し、すべての課題の共通入力とせよ。

\begin{lstlisting}[numbers=none, xleftmargin=2\zw, title={events.csv}]
detector,time_ns,signal,valid
A,1000,12.4,true
B,1180,9.8,true
A,1300,15.1,true
B,1510,,false
A,1760,11.0,true
B,2050,13.2,true
\end{lstlisting}

\begin{enumerate}
\item \textbf{読み込み直後の表の検査}。
\code{pd.read\_csv()} で読み、\code{shape}、\code{columns}、\code{dtypes}、\code{head(3)}、\code{isna().sum()} を表示せよ。行数は 6、列数は 4、\code{signal} の欠損数は 1 になる。\code{head(3)} の \code{3} と、\code{isna().sum()} の \code{sum()} がどの方向に集計しているかを説明せよ。

\item \textbf{条件抽出と列追加}。
\code{valid} が真で \code{signal >= 12.0} の行を \code{loc} で抽出し、\code{signal / signal.mean()} を \code{signal\_ratio} 列として追加せよ。抽出行は 3 件になる。複数条件を括弧で囲み、\code{and} ではなく \code{\&} を使う理由を説明せよ。

\item \textbf{検出器別集計と別表の結合}。
\code{groupby("detector")} と \code{agg()} を使い、有効な信号の件数、平均、標準偏差を求めよ。平均は A が約 12.8333、B が 11.5 になる。さらに、A と B の較正係数を持つ DataFrame を作り、\code{merge(..., on="detector", how="left")} で結合せよ。\code{on} と \code{how} の役割を説明せよ。

\item \textbf{一括読み込みと分割読み込みの比較}。
本文の \code{chunksize=2} の方法で件数と平均を求め、一括読み込みから得た結果と \code{pd.testing.assert\_frame\_equal()} または \code{np.allclose()} で一致を確認せよ。チャンク平均の単純平均が一般には誤る理由を、件数が異なる 2 チャンクの例を使って説明せよ。

\item \textbf{再利用可能な形式への保存}。
集計表を CSV と Parquet の両方へ保存し、読み戻して列名と値を確認せよ。\code{index=False} を付けた場合に何が保存されないかを説明し、Parquet の読み込みに追加ライブラリが必要な環境では、表示された例外名も記録せよ。
\end{enumerate}

pandas で表処理の基本と結果の検証方法を確認した後は、同じ処理を別の DataFrame 実装で書くと、設計上の違いが明確になる。次章では、Polars による index 非依存の表処理、列式、遅延実行を扱う。
